\documentclass[11pt]{article}
\usepackage[margin=1in]{geometry}
\usepackage{amsmath,amssymb,amsthm}
\usepackage{booktabs}
\usepackage{array}
\usepackage{graphicx}
\usepackage{hyperref}
\hypersetup{hidelinks}

\newtheorem{theorem}{Theorem}
\newtheorem{lemma}[theorem]{Lemma}
\newtheorem{corollary}[theorem]{Corollary}
\newtheorem{proposition}{Proposition}
\theoremstyle{definition}
\newtheorem{definition}{Definition}
\newtheorem*{remark}{Remark}

\setlength{\parskip}{4pt plus 1pt minus 1pt}
\setlength{\parindent}{0pt}

\title{The Cost of Computation:\\Shared Molecular Machinery Between Learning and Cancer}
\author{Patrick D.\ McCarthy}
\date{}

\begin{document}
\maketitle

\begin{abstract}
Paper~3 of this series proved that knowledge and action are
inseparable under bounded context: the machinery that stores
knowledge is the same machinery that executes action. This paper
identifies the physical substrate of that inseparability.

Neurons break their own DNA to learn. Topoisomerase~II$\beta$
creates double-strand breaks in promoter regions during memory
formation \cite{Madabhushi2015}. The chromatin remodeling genes
that repair this damage---\emph{ARID1A}, \emph{KMT2C},
\emph{SMARCA4}, \emph{DNMT3A}---are the same genes whose somatic
mutation drives cancer progression. This is not coincidence. It is
the molecular consequence of K/A inseparability: the read/write
machinery for knowledge is the same machinery whose failure
enables cancer.

We prove three claims. First, learning causes targeted DNA
double-strand breaks (established biology). Second, the genes that
repair these breaks overlap with cancer driver genes (checkable
from public data). Third, this overlap is a structural consequence
of bounded context: any system that modifies its own substrate to
store knowledge will share failure modes with uncontrolled
substrate modification. We derive predictions and state what we do
not yet know.
\end{abstract}

% ------------------------------------------------------------------
\section{Introduction}
\label{sec:intro}
% ------------------------------------------------------------------

Learning has a physical cost. Neurons break their own DNA to learn.

Topoisomerase II$\beta$ creates double-strand breaks in the promoter
regions of immediate early genes during memory formation
\cite{Madabhushi2015}. These breaks are necessary; without them,
the chromatin cannot open fast enough for rapid transcriptional
response. The cell repairs the breaks, but imperfectly.
Each learning event leaves a residue of somatic mutation.

This is not damage from an external insult. It is the cost of
computation. The cell breaks its DNA on purpose, reads the gene,
repairs the break, and moves on. Over a lifetime, the repair
errors accumulate.

The question is: does the \emph{efficiency} of learning determine
the rate of accumulation?

We argue that it does.

\subsection{The cost per bit}

In the bounded-context framework \cite{McCarthy2026a}, an agent
reduces its uncertainty about the world $H(W \mid K)$ by acquiring
knowledge $K$ through a channel with capacity $C_n$. The number
of computational cycles required to reduce uncertainty by one bit
is bounded below by $1/C_n$. Each cycle, in a biological neural
network, involves gene expression changes that require DNA breaks.

The mutation cost per bit of uncertainty reduced is therefore:

\begin{equation}
\label{eq:cost}
\mu = \frac{\delta \cdot N_{\text{cycles}}}{\Delta H}
    \geq \frac{\delta}{C_n}
\end{equation}

where $\delta$ is the per-cycle probability of a repair error
(approximately $10^{-9}$ per base pair per repair event
\cite{Lindahl1993}), $N_{\text{cycles}}$ is the number of
computation cycles, and $\Delta H$ is the information gained.

An efficient learner has high $C_n$: each cycle reduces uncertainty
productively. The DNA damage cost per bit is near the minimum
$\delta / C_n$.

Any process that reduces effective $C_n$ without reducing
$N_{\text{cycles}}$ increases the mutation cost per bit.
The system computes at the same rate but learns less per cycle.
Over decades, the difference compounds.

\subsection{Connection to Paper~3}

Paper~3 \cite{McCarthy2026c} proved that knowledge and action
are inseparable under bounded context: the same substrate that
stores propositions must also execute inference over them.
In biological systems, that substrate is chromatin. The genes
that control which other genes are readable---the chromatin
remodeling complex---are simultaneously the knowledge store
(their configuration determines what the cell ``knows'') and
the action machinery (their activity determines what the cell
can do).

This paper identifies the consequence: because the knowledge
machinery and the action machinery are the same physical genes,
damage from learning (knowledge) and damage enabling cancer
(uncontrolled action) share a molecular target.

% ------------------------------------------------------------------
\section{The Shared Molecular Machinery}
\label{sec:machinery}
% ------------------------------------------------------------------

The bounded-context framework predicts that learning and cancer
should share molecular machinery. They do.

\subsection{Chromatin remodeling: the read/write heads}

The genes that control chromatin accessibility---determining which
genes are readable in which cellular context---are mutated in both
cognitive disorders and cancer:

\begin{itemize}
\item \emph{ARID1A}: mutated in 6--8\% of all cancers; germline
variants cause intellectual disability \cite{Santen2012}.
\item \emph{KMT2C}/\emph{KMT2D}: among the most frequently mutated
genes in cancer; germline \emph{KMT2D} loss causes Kabuki syndrome
(intellectual disability) \cite{Miyake2013}.
\item \emph{SMARCA4}: mutated in lung, ovarian, and other cancers;
germline loss causes Coffin-Siris syndrome (intellectual
disability) \cite{Tsurusaki2012}.
\item \emph{DNMT3A}: mutated in AML; germline variants cause
Tatton-Brown-Rahman syndrome (overgrowth and intellectual
disability) \cite{Tatton-Brown2014}.
\end{itemize}

These are not coincidences. In the bounded-context framework,
chromatin remodeling genes are the molecular implementation of
context switching. They decide which genes are accessible
(loaded into the context window) in a given cell state. When they
fail by somatic mutation, the cell loses context control: it
expresses genes that do not belong in its current state. That is
cancer. When they fail by germline mutation, the organism loses
context control at the neural level: it cannot efficiently switch
cognitive contexts. That is intellectual disability.

\subsection{TP53: the context integrity checkpoint}

TP53 is the most frequently mutated gene in cancer
\cite{Kandoth2013} and a central node in the coupling-channel
framework \cite{McCarthy2026e}. It functions as a context
integrity check: if a cell's DNA is too damaged to maintain
faithful gene expression context, p53 activates apoptosis.

In our framework, p53 is the kill switch for context window
corruption. Its dual role in cancer (loss enables tumor growth)
and neurodegeneration (activation kills neurons) is not a paradox.
It is the same function---context integrity enforcement---applied
in two domains where the optimal threshold differs.

The inverse comorbidity between cancer and Alzheimer's
\cite{Driver2012, Musicco2013} reflects this: how aggressively an
organism enforces context integrity determines which failure mode
it is susceptible to. Cancer requires lax enforcement (corrupted
cells survive). Neurodegeneration requires strict enforcement
(compromised neurons are killed). The same axis, different
thresholds \cite{Behrens2009}.

\subsection{The immune system: knowledge stored in DNA}

The immune system stores knowledge through somatic hypermutation
in B cells. This is a literal implementation of the bounded-context
framework:

\begin{enumerate}
\item Encounter an antigen (new evidence).
\item Mutate DNA to generate antibody variants (context switching
at the molecular level).
\item Select the variant that best matches (uncertainty reduction).
\item Store as memory B cells (knowledge encoded in DNA).
\end{enumerate}

Autoimmune diseases are context indexing failures: the retrieval
system misidentifies self as threat, loading the wrong context.
Autoimmune diseases are associated with elevated cancer risk
\cite{Franks2012}---consistent with our framework: a system that
loads wrong targets imposes DNA damage costs on unintended cells.

% ------------------------------------------------------------------
\section{The Epidemiological Signal}
\label{sec:signal}
% ------------------------------------------------------------------

If shared machinery is real, psychiatric conditions should show
altered cancer risk. The epidemiological literature contains
a striking pattern: some conditions elevate risk, others reduce it.

\begin{table}[ht]
\centering
\caption{Observed cancer risk in psychiatric and neurological
conditions. Direction from meta-analyses and large registry
studies. We report what is observed, not what we can explain.}
\label{tab:signal}
\small
\begin{tabular}{@{}llll@{}}
\toprule
\textbf{Condition} & \textbf{Cancer risk} &
\textbf{Key reference} & \textbf{Status} \\
\midrule
Depression & $\uparrow$ & \cite{Chida2008} & Replicated \\
Insomnia & $\uparrow$ & \cite{Lu2017} & Replicated \\
Autoimmune & $\uparrow$ & \cite{Franks2012} & Replicated \\
TBI & $\uparrow$ (glioma) & \cite{Chen2012tbi} & Replicated \\
OCD & $\uparrow$ (mild) & \cite{Shen2013} & Single study \\
Anxiety (GAD) & $\uparrow$ (mild) & \cite{Batty2017} & Mixed \\
Bipolar & $\uparrow$ (mild) & \cite{McGinty2012} & Mixed \\
\midrule
Schizophrenia & $\downarrow$ & \cite{Catts2008, Ji2013} & Replicated \\
Alzheimer's & $\downarrow$ & \cite{Driver2012, Musicco2013} & Replicated \\
Parkinson's & $\downarrow$ (most solid) & \cite{Bajaj2010} & Replicated \\
Down syndrome & $\downarrow$ (solid) & \cite{Hasle2000} & Replicated \\
\midrule
PTSD & null after adjustment & \cite{Gradus2015} & Open question \\
Autism & mixed & \cite{Darbro2016} & Heterogeneous \\
\bottomrule
\end{tabular}
\end{table}

\subsection{What we observe}

The table divides. Conditions with intact but inefficient
computation show elevated risk. Conditions where the computation
or context-switching machinery is disabled show reduced risk.

The bounded-context framework offers a structural explanation:
conditions that overwork the chromatin remodeling machinery
(rumination, replay, inefficient context switching) accumulate
somatic damage in those genes. Conditions that disable the
machinery (schizophrenia, severe neurodegeneration) cannot
accumulate the specific somatic mutations that cancer requires---and
cancer cells themselves need functional chromatin remodeling
to execute the gene expression reprogramming necessary for
tumor progression.

\subsection{What we cannot yet sign}

PTSD is predicted by our framework to elevate cancer risk
through the context-lock mechanism: traumatic context replays
involuntarily, each replay incurring DNA damage cost with
zero information gain ($\Delta H = 0$). But the largest
epidemiological study (Gradus et al.\ 2015, $n > 40{,}000$
Danish veterans) found no association after adjustment for
behavioral confounders. The signal may be real but confounded,
or our classification may be wrong. This is an open test.

Autism is heterogeneous. Mild rigidity may elevate cost-per-bit
(the system over-processes). Severe rigidity may reduce
context-switching frequency (fewer cycles, less DNA damage).
The mixed epidemiological evidence \cite{Darbro2016} is
consistent with a spectrum, but we cannot sign the relationship
without stratification by severity.

We do not have a taxonomy of context failure modes. The original
version of this paper proposed sixteen classifications. That was
wrong---we do not know whether the correct number is seven or
seventy. What we know is that the molecular machinery is shared,
the epidemiological signal divides along a recognizable axis,
and the framework makes predictions about which direction each
condition should go. Some of those predictions are untested.

% ------------------------------------------------------------------
\section{Prior Work}
\label{sec:prior}
% ------------------------------------------------------------------

\textbf{Activity-dependent DNA damage.} Madabhushi et al.\
\cite{Madabhushi2015} demonstrated that neuronal activity causes
targeted DNA double-strand breaks. Subsequent work showed that
these breaks accumulate with age and correlate with
neurodegenerative disease \cite{Thadathil2021}. We extend this to
cancer risk: if the break rate depends on computational efficiency,
then conditions that reduce efficiency increase the break-to-knowledge
ratio.

\textbf{Psychoneuroimmunology.} The field of
psychoneuroimmunology has studied how psychological states
affect immune function and cancer outcomes \cite{Kiecolt-Glaser2002}.
Our contribution is mechanistic specificity: not ``stress causes
cancer through immune suppression'' but ``context-switching waste
causes somatic mutation in the specific genes that control chromatin
accessibility, through the specific mechanism of topoisomerase-mediated
DNA breaks.''

\textbf{Chromatin remodeling and schizophrenia.} Singh et al.\
\cite{Singh2016} showed that rare schizophrenia risk variants are
enriched in chromatin remodeling genes. If these genes are already
non-functional from germline variants, they cannot be further
corrupted by somatic mutation to enable cancer---and cancer cells
require functional chromatin remodeling to progress. This explains
the reduced cancer risk in schizophrenia without invoking
lifestyle factors or antipsychotic effects.

\textbf{CHIP as evidence of chromatin vulnerability.}
Clonal hematopoiesis of indeterminate potential (CHIP) is
driven predominantly by mutations in \emph{DNMT3A}, \emph{TET2},
and \emph{ASXL1}---all chromatin remodeling genes, accounting
for approximately 80\% of CHIP mutations \cite{Jaiswal2014,
Steensma2015}. This concentration is not predicted by uniform
mutagenesis models. It is consistent with our framework:
chromatin remodeling genes are under elevated mutational pressure
because they are the most active read/write machinery in the genome.

\textbf{Coupling channels and cancer.} In the companion paper
\cite{McCarthy2026e}, we showed that the number of coupling
channels disrupted predicts cancer survival across 73,593 patients.
The coupling channels include chromatin remodeling and DNA
methylation---precisely the channels this paper identifies as
the shared molecular substrate between cognitive computation
and cancer.

\textbf{Bounded context and intelligence.} In \cite{McCarthy2026d},
we derived the structural invariants of intelligence from bounded
context. The present paper adds a biological cost to evaluation-driven
descent: each step of the intelligence loop has a DNA damage cost.
Intelligence is not free. Its cost is somatic mutation.

% ------------------------------------------------------------------
\section{Predictions}
\label{sec:predictions}
% ------------------------------------------------------------------

The theory makes specific, falsifiable predictions. Most were
derived before we knew whether the epidemiological data supported
them. The ones that held (schizophrenia $\downarrow$, Alzheimer's
$\downarrow$, depression $\uparrow$) were not fitted to the data.
They follow from the cost-per-bit equation and the shared machinery.
The untested ones are genuine forecasts.

\subsection{Predictions from the cost-per-bit equation}

\begin{enumerate}
\item \textbf{Chromatin specificity.} The elevated somatic
mutation burden in conditions with inefficient context switching
is concentrated in chromatin remodeling genes, not uniformly
distributed across the genome. Testable: compare mutation spectra
in CHIP-positive individuals with and without psychiatric history.

\item \textbf{Efficiency, not distress, determines risk.}
Interventions that improve context efficiency (bits per cycle)
are predicted to be protective. Interventions that reduce distress
without improving efficiency are not. Specifically:
stimulant treatment for ADHD (improves efficiency, predicted
protective \cite{Dalsgaard2018}), SSRIs that reduce rumination
(improves efficiency, some evidence \cite{Chubak2011}),
benzodiazepines (reduce distress but impair formation, predicted
not protective \cite{Kao2017}). The discriminating test:
matched cohorts on the same anxiolytic vs.\ cognitive therapy,
comparing chromatin remodeling VAF from cfDNA.

\item \textbf{Antipsychotics restore cancer risk.} Antipsychotic
treatment in schizophrenia should increase cancer risk relative to
unmedicated patients by restoring chromatin remodeling function.
Some evidence exists \cite{Solmi2020}. The prediction: the effect
should correlate with degree of cognitive improvement, not with
metabolic side effects.

\item \textbf{PTSD direction.} If PTSD is truly a context-lock
disorder (replay with $\Delta H = 0$), then cancer risk should be
elevated in PTSD patients with high flashback frequency,
specifically through chromatin remodeling mutations, after
controlling for behavioral confounders. The null result in
Gradus et al.\ \cite{Gradus2015} may reflect insufficient
stratification by symptom severity. This is the prediction we
are least confident about.

\item \textbf{Non-carcinogen addictions carry cancer risk.}
Behavioral addictions (gambling, gaming) without chemical
carcinogen exposure should show elevated somatic mutation burden
through the context-lock mechanism alone.
\end{enumerate}

\subsection{Predictions from the shared machinery}

\begin{enumerate}
\setcounter{enumi}{5}
\item \textbf{Chromatin remodeling gene overlap is checkable.}
The genes whose germline loss causes intellectual disability
and whose somatic mutation drives cancer should overlap at a rate
significantly above chance. This is directly testable from
OMIM + COSMIC + ClinVar data. If the overlap is not significant,
the framework has a problem.

\item \textbf{CHIP onset correlates with cognitive load.}
If context-switching drives chromatin mutation, then CHIP onset
should correlate with lifetime cognitive load after controlling
for age and carcinogen exposure. High-information-processing
occupations should show earlier CHIP than low-switching occupations.

\item \textbf{Early-onset cancers are enriched for chromatin
mutations.} The cancers driving the early-onset trend (colorectal,
breast, pancreatic in adults under 50 \cite{Sung2024}) should show
higher rates of chromatin remodeling mutations than late-onset
counterparts, after controlling for stage and subtype.
\end{enumerate}

\subsection{The computer analogy as prediction}

A computer's CPU generates heat proportional to its clock speed,
independent of whether the computation is useful. A CPU running
an infinite loop at full speed generates the same heat as one
solving a real problem. The damage (thermal wear, electromigration)
accumulates either way.

Neurons are the same. The DNA damage cost is proportional to
computation cycles, not to useful bits learned. A brain in
rumination, in traumatic replay, in anxious scanning---runs at
full clock speed, generating the same molecular wear as productive
learning. The cost-per-bit equation formalizes this: when $\Delta H
\to 0$ but $N_{\text{cycles}}$ remains constant, $\mu \to \infty$.

This analogy is suggestive of a broader prediction: any biological
system that modifies its own substrate to store information will
exhibit a cancer-like failure mode when that modification becomes
inefficient. The immune system's somatic hypermutation already
provides one example (lymphoma as the failure mode of adaptive
immunity). The framework predicts others exist.

\subsection{What these predictions test}

If most of these predictions hold, the framework is confirmed
independently of the data it was built on. If the chromatin
overlap (Prediction~6) fails, the molecular basis is wrong.
If the efficiency distinction (Prediction~2) fails, the
cost-per-bit mechanism is wrong. If the PTSD prediction
(Prediction~4) fails, our classification of context failure
modes needs revision---but the core machinery claim survives.

We state what we do not know. We do not know the correct
taxonomy of context failure modes. We do not know whether the
inverted-U relationship between severity and risk is real or an
artifact of how we classify severity. We do not know the
quantitative parameters ($\delta$, $C_n^{\text{eff}}$ per
condition). What we know is that the machinery is shared,
the signal exists, and the predictions are falsifiable.

% ------------------------------------------------------------------
\section{Discussion}
\label{sec:discussion}
% ------------------------------------------------------------------

The central claim of this paper is simple. Learning breaks DNA.
The genes that repair the breaks are the same genes whose failure
enables cancer. This is a structural consequence of K/A
inseparability: any system where the knowledge substrate and the
action substrate are physically identical will share failure modes
between knowledge acquisition and uncontrolled action.

This creates a direct, testable, molecular link between the
efficiency of cognition and the risk of cancer---not through stress
hormones, not through behavioral factors, not through immune
suppression, but through the physical cost of computation in a
biological substrate.

The deepest implication is for the relationship between knowledge
and biology. If learning has a DNA damage cost specific to the
genes that control gene expression context, then what we know is
physically inscribed in the wear pattern of our chromatin
remodeling machinery. Not in the sequence itself (like the immune
system's somatic hypermutation), but in the error pattern on the
read/write heads.

% ------------------------------------------------------------------
\section*{Conflict of Interest}
The author declares no competing interests.

\section*{Funding}
This research received no external funding.

% ------------------------------------------------------------------
% References
% ------------------------------------------------------------------

\begin{thebibliography}{99}

\bibitem{McCarthy2026a}
P.~D.~McCarthy, ``Graph structure is necessary for information
preservation under bounded context,'' 2026.

\bibitem{McCarthy2026c}
P.~D.~McCarthy, ``Knowledge as normalization: the inseparability of
knowledge and action under bounded context,'' 2026.

\bibitem{McCarthy2026d}
P.~D.~McCarthy, ``Evaluation-driven descent, encoding permanence,
and the structural invariants of intelligence,'' 2026.

\bibitem{McCarthy2026e}
P.~D.~McCarthy, ``Coupling-channel count predicts cancer survival
better than mutation count,'' 2026.

\bibitem{Madabhushi2015}
R.~Madabhushi et~al., ``Activity-induced {DNA} breaks govern the
expression of neuronal early-response genes,'' \emph{Cell},
vol.~161, pp.~1592--1605, 2015.

\bibitem{Lindahl1993}
T.~Lindahl, ``Instability and decay of the primary structure of
{DNA},'' \emph{Nature}, vol.~362, pp.~709--715, 1993.

\bibitem{Chen2012tbi}
Y.-H.~Chen et~al., ``Increased risk of brain cancer in people
with traumatic brain injury,'' \emph{Br.~J.~Cancer}, 2012.

\bibitem{Shen2013}
C.-C.~Shen et~al., ``Risk of cancer in patients with obsessive-compulsive
disorder,'' \emph{J.~Clin.~Psychiatry}, 2013.

\bibitem{Chida2008}
Y.~Chida et~al., ``Does psychosocial well-being predict cancer
incidence?'' \emph{Nature Rev.~Clin.~Oncol.}, vol.~5,
pp.~466--475, 2008.

\bibitem{Batty2017}
G.~D.~Batty et~al., ``Psychosocial factors and cancer incidence,''
\emph{BMJ}, 2017.

\bibitem{McGinty2012}
E.~E.~McGinty et~al., ``Cancer incidence in a sample of
{M}aryland residents with serious mental illness,''
\emph{Psychiatric Services}, vol.~63, pp.~714--717, 2012.

\bibitem{Catts2008}
V.~S.~Catts et~al., ``Cancer incidence in patients with
schizophrenia,'' \emph{Schizophrenia Research}, vol.~102,
pp.~137--162, 2008.

\bibitem{Ji2013}
J.~Ji et~al., ``Cancer risk in schizophrenia,'' \emph{BMC Cancer},
vol.~13, 2013.

\bibitem{Singh2016}
T.~Singh et~al., ``Rare schizophrenia risk variants are enriched
in chromatin remodeling genes,'' \emph{Nature Neuroscience},
vol.~19, pp.~571--577, 2016.

\bibitem{Solmi2020}
M.~Solmi et~al., ``Cancer risk in antipsychotic-treated
schizophrenia,'' \emph{Lancet Psychiatry}, 2020.

\bibitem{Driver2012}
J.~A.~Driver et~al., ``Inverse association between cancer and
{A}lzheimer's disease,'' \emph{BMJ}, vol.~344, e1442, 2012.

\bibitem{Musicco2013}
M.~Musicco et~al., ``Inverse occurrence of cancer and {A}lzheimer
disease,'' \emph{Neurology}, vol.~81, pp.~322--328, 2013.

\bibitem{Behrens2009}
M.~I.~Behrens et~al., ``Inverse relationship of {A}lzheimer's
disease and cancer,'' \emph{JAMA Neurol.}, 2009.

\bibitem{Darbro2016}
B.~W.~Darbro et~al., ``Autism linked to increased oncogene
mutations but decreased cancer rate,'' \emph{PLOS ONE}, vol.~11,
e0149041, 2016.

\bibitem{Lu2017}
Y.~Lu et~al., ``Insomnia and cancer risk,'' \emph{Sleep Medicine
Reviews}, 2017.

\bibitem{Bajaj2010}
A.~Bajaj et~al., ``{P}arkinson's disease and cancer risk,''
\emph{Cancer Causes Control}, vol.~21, pp.~697--707, 2010.

\bibitem{Hasle2000}
H.~Hasle et~al., ``Cancer incidence in {D}own syndrome,''
\emph{Lancet}, vol.~355, pp.~165--169, 2000.

\bibitem{Santen2012}
G.~W.~E.~Santen et~al., ``Mutations in \emph{ARID1A} cause
intellectual disability,'' \emph{Am.~J.~Hum.~Genet.}, 2012.

\bibitem{Miyake2013}
N.~Miyake et~al., ``\emph{KMT2D} and {K}abuki syndrome,''
\emph{Am.~J.~Med.~Genet.}, 2013.

\bibitem{Tsurusaki2012}
Y.~Tsurusaki et~al., ``Mutations in \emph{SMARCA4} cause
{C}offin-{S}iris syndrome,'' \emph{Nature Genetics}, 2012.

\bibitem{Tatton-Brown2014}
K.~Tatton-Brown et~al., ``Mutations in \emph{DNMT3A} cause
{T}atton-{B}rown-{R}ahman syndrome,'' \emph{Am.~J.~Hum.~Genet.},
2014.

\bibitem{Kandoth2013}
C.~Kandoth et~al., ``Mutational landscape and significance across
12 major cancer types,'' \emph{Nature}, vol.~502, pp.~333--339,
2013.

\bibitem{Kiecolt-Glaser2002}
J.~K.~Kiecolt-Glaser et~al., ``Psychoneuroimmunology and cancer,''
\emph{Ann.~Oncol.}, 2002.

\bibitem{Thadathil2021}
N.~Thadathil et~al., ``{DNA} double-strand breaks accumulate in
aging brain,'' \emph{Frontiers in Aging Neuroscience}, 2021.

\bibitem{Franks2012}
A.~L.~Franks and J.~E.~Slansky, ``Autoimmune disease and cancer,''
\emph{Autoimmunity}, 2012.

\bibitem{Jaiswal2014}
S.~Jaiswal et~al., ``Age-related clonal hematopoiesis associated
with adverse outcomes,'' \emph{N.~Engl.~J.~Med.}, vol.~371,
pp.~2488--2498, 2014.

\bibitem{Steensma2015}
D.~P.~Steensma et~al., ``Clonal hematopoiesis of indeterminate
potential and its distinction from myelodysplastic syndromes,''
\emph{Blood}, vol.~126, pp.~9--16, 2015.

\bibitem{Gradus2015}
J.~L.~Gradus et~al., ``Posttraumatic stress disorder and cancer
risk,'' \emph{JAMA Psychiatry}, vol.~72, pp.~1179--1185, 2015.

\bibitem{Chubak2011}
J.~Chubak et~al., ``Antidepressant use and colorectal cancer
risk,'' \emph{Cancer Causes Control}, 2011.

\bibitem{Dalsgaard2018}
S.~Dalsgaard et~al., ``{ADHD} medication and cancer risk,''
\emph{J.~Child Psychol.~Psychiatry}, 2018.

\bibitem{Kao2017}
C.-H.~Kao et~al., ``Benzodiazepine use and risk of cancer:
a population-based cohort study,'' \emph{Psycho-Oncology},
2017.

\bibitem{Sung2024}
H.~Sung et~al., ``Global trends in incidence of early-onset
cancer,'' \emph{BMJ Oncology}, 2024.

\end{thebibliography}

\end{document}
